\documentclass[12pt,twoside]{article}

\input{macros-sp26}
\newcommand{\thelecturenum}{7}

\title{CS336 Lecture 7}

\begin{document}

\handout{Lecture \thelecturenum: Systems Thinking}

\setlength{\parindent}{0pt}

\subsection*{Scaling the Frontier Bottleneck}

\textbf{Problem:} Compute (ALUs) is far from inputs/outputs (data). 

\begin{itemize}
    \item Last week: parallelism within a single CPU or GPU.
    \item This week: parallelism across multiple CPUs or GPUs.
\end{itemize}

\textbf{Solution:} Orchestrate computation to avoid data transfer bottlenecks.

\begin{itemize}
    \item Single node, single CPU-GPU: L1 cache / shared memory (fastest).
    \item Single node, single CPU-GPU: HBM ("fast").
    \item Single node, multi-GPU: NVLink / NVSwitch (slow).
    \item Multi-node, multi-GPU: Infiniband / Ethernet (slowest).
\end{itemize}

\textbf{Mechanics:} How to maximize compute and minimize latency?

\begin{itemize}
    \item Last week: reduce memory accesses via fusion/tiling.
    \item This week: reduce communication across GPUs/nodes via replication/sharding.
\end{itemize}

\textbf{Question:} Why do multi-GPU?

\begin{enumerate}
    \item Your parameters (optimizer state + gradients + activations) don't fit on a single GPU.
    \item You want to use more GPUs (more FLOPs) to train faster.
\end{enumerate}

\subsection*{Distributing the Computation}

\textbf{Collective operations} are the conceptual primitives used for distributed programming:

\begin{itemize}
    \item These are classic in the parallel programming literature from the 1980s.
    \item \textit{Collective} means that you specify a general communication pattern across many devices.
    \item This can be better/faster than managing point-to-point communication yourself.
\end{itemize}

Setup for distributed computation and communication:

\begin{itemize}
    \item \textbf{Rank}: a particular CPU or GPU (e.g., 0, 1, 2, 3).
    \item \textbf{World size}: total number of devices (e.g., 4).
\end{itemize}

Operations for distributed systems and networks:

\begin{itemize}
    \item Broadcast, scatter, gather, reduce (foundations).
    \item All-gather, reduce-scatter, all-reduce (workhorse).
    \item All-to-all (for MoEs).
\end{itemize}

\textbf{Broadcast}: copy from rank 0 to all ranks.

\begin{minted}{python}
    # Input
    rank0 = tensor([0., 1, 2, 3])
    
    # Output
    rank0 = tensor([0., 1, 2, 3])
    rank1 = tensor([0., 1, 2, 3])
    rank2 = tensor([0., 1, 2, 3])
    rank3 = tensor([0., 1, 2, 3])
\end{minted}

\textbf{Scatter} tensor on rank 0 to all ranks.

\begin{minted}{python}
    # Input
    rank0 = tensor([0., 1, 2, 3])
    
    # Output
    rank0 = tensor([0.])
    rank1 = tensor([1.])
    rank2 = tensor([2.])
    rank3 = tensor([3.])
\end{minted}

\textbf{Gather} pieces from all ranks to rank 0 (opposite of scatter).

\begin{minted}{python}
    # Input
    rank0 = tensor([0.])
    rank1 = tensor([1.])
    rank2 = tensor([2.])
    rank3 = tensor([3.])
    
    # Output
    rank0 = tensor([0., 1, 2, 3])
\end{minted}

\textbf{Reduce} pieces from all ranks to rank 0, applying some operation (e.g, sum, min, max).

\begin{minted}{python}
    # Input
    rank0 = tensor([0.])
    rank1 = tensor([1.])
    rank2 = tensor([2.])
    rank3 = tensor([3.])
    
    # Output
    rank0 = tensor([6.])  # Sum of all ranks (0 + 1 + 2 + 3)
\end{minted}

\textbf{All-gather}: perform gather to all ranks, not just rank 0.

\begin{minted}{python}
    # Input
    rank0 = tensor([0.])
    rank1 = tensor([1.])
    rank2 = tensor([2.])
    rank3 = tensor([3.])
    
    # Output
    rank0 = tensor([0., 1, 2, 3])
    rank1 = tensor([0., 1, 2, 3])
    rank2 = tensor([0., 1, 2, 3])
    rank3 = tensor([0., 1, 2, 3])
\end{minted}

\textbf{Reduce-scatter}: perform reduce on each dimension, scatter results.

\begin{minted}{python}
    # Input
    rank0 = tensor([0., 1, 2, 3])
    rank1 = tensor([1., 2, 3, 4])
    rank2 = tensor([2., 3, 4, 5])
    rank3 = tensor([3., 4, 5, 6])
    
    # Output
    rank0 = tensor([6.])  # Sum along dim 0 (0 + 1 + 2 + 3)
    rank1 = tensor([10.]) # Sum along dim 1 (1 + 2 + 3 + 4)
    rank2 = tensor([14.]) # Sum along dim 2 (2 + 3 + 4 + 5)
    rank3 = tensor([18.]) # Sum along dim 3 (3 + 4 + 5 + 6)
\end{minted}

\textbf{All-reduce} \(=\) reduce-scatter \(+\) all-gather.

\begin{minted}{python}
    # Input
    rank0 = tensor([0., 1, 2, 3])
    rank1 = tensor([1., 2, 3, 4])
    rank2 = tensor([2., 3, 4, 5])
    rank3 = tensor([3., 4, 5, 6])
    
    # Output
    rank0 = tensor([6., 10, 14, 18])
    rank1 = tensor([6., 10, 14, 18])
    rank2 = tensor([6., 10, 14, 18])
    rank3 = tensor([6., 10, 14, 18])
\end{minted}

\textbf{All-to-all}: each rank sends each other rank some tensor (most general).

\begin{minted}{python}
    # Input
    rank0 = tensor([0., 1, 2, 3])    
    rank1 = tensor([4., 5, 6, 7])
    rank2 = tensor([8., 9, 10, 11])
    rank3 = tensor([12., 13, 14, 15])
    
    # Output
    rank0 = tensor([0, 4, 8, 12])
    rank1 = tensor([1, 5, 9, 13])
    rank2 = tensor([2, 6, 10, 14])
    rank3 = tensor([3, 7, 11, 15])
\end{minted}

Notes: 

\begin{itemize}
    \item Useful for MoEs: each rank has split of data and subset of experts; need to route data to experts.
    \item For balanced splits, all-to-all looks like transpose.
    \item Also handles unbalanced splits (but want splits to be as balanced as possible).
\end{itemize}

Way to remember the terminology:

\begin{itemize}
    \item Reduce - performs some associative/commutative operation (sum, min, max).
    \item Scatter is inverse of gather.
    \item All - means destination is all devices.
\end{itemize}

\subsection*{Hardware Architecture}

Classic (in the home):

\begin{itemize}
    \item GPUs on same node communicate via a PCI(e) bus (v7.0, 16 lanes \(\implies\) 242 GB/s).
    \item GPUs on different nodes communicate via Ethernet (200 MB/s).
\end{itemize}

Modern (in the data center):

\begin{itemize}
    \item 8 GPUs per node, connected by NVLink to an NVSwitch (B200s' NVLink 5.0 gets 1.8 TB/s; HBM was 8 TB/s).
    \item 256 nodes per pod, connected by Infiniband (PCIe \(\to\) HCA / Infiniband NIC \(\to\) Infiniband cable) (0.05 TB/s).
    \item N pods per cluster / datacenter, connected by Ethernet (PCIe \(\to\) CPU).
\end{itemize}

Bypassing the CPU:

\begin{itemize}
    \item Ethernet requires passing through the CPU (copying data to kernel socket buffer, build TCP packets, copy to NIC ring buffer).
    \item Remote Direct Memory Access (RDMA) - allows one GPU to directly read/write another GPU's memory without involving the CPU.
    \item Infiniband supports RDMA, but standard Ethernet does not.
\end{itemize}

Advancements:

\begin{itemize}
    \item GB200/GB300 NVL72 - 8 GPUs per tray, 9 trays per rack \(\to\) 72 GPUs in one NVLink domain.
    \item RDMA over Converged Ethernet (RoCE) - Ethernet bypasses CPU, similar but cheaper/weaker than Infiniband, used by Meta.
\end{itemize}

\subsubsection*{NVIDIA Collective Communication Library (NCCL)}

NCCL translates collective operations into low-level packets that are sent between GPUs.

\begin{itemize}
    \item Detects topology of hardware (e.g, number of nodes, switches, NVLink/PCIe).
    \item Optimizes the path between GPUs.
    \item Launches GPU kernels to send/receive data.
\end{itemize}

\subsection*{Torch Distributed Operations}

PyTorch distributed library (\textbf{torch.distributed}).

\begin{itemize}
    \item Provides clean interface for collective operations (e.g., all\_gather\_into\_tensor).
    \item Supports multiple backends for different hardware: gloo (CPU), nccl (GPU).
    \item Also supports higher-level algorithms (e.g, FullyShardedDataParallel) [not used in this course].
\end{itemize}

\begin{minted}{python}
def collective_operations_main(rank: int, world_size: int):  
    """This function is running asynchronously 
       for each process (rank = 0, ..., world_size - 1)."""
    setup(rank, world_size)
    
    ### All-reduce (dist = torch.distributed)
    dist.barrier() 
    data = tensor([0., 1, 2, 3], device=cuda(rank)) + rank
    print(f"Rank {rank} [before all-reduce]: {data}", flush=True)
    dist.all_reduce(tensor=data, 
                    op=dist.ReduceOp.SUM, async_op=False) 
    print(f"Rank {rank} [after all-reduce]: {data}", flush=True)
    
    ### Reduce-scatter
    dist.barrier()
    input = torch.arange(world_size, 
                         dtype=torch.float32, 
                         device=cuda(rank)) + rank
    output = torch.empty(1, device=cuda(rank))
    print(f"Rank {rank} [before reduce-scatter]: input = {input}, 
          output = {output}", flush=True)
    dist.reduce_scatter_tensor(output=output, 
                               input=input, 
                               op=dist.ReduceOp.SUM, 
                               async_op=False)
    print(f"Rank {rank} [after reduce-scatter]: input = {input}, 
          output = {output}", flush=True)
    
    ### All-gather
    dist.barrier()
    input = output
    output = torch.empty(world_size, device=cuda(rank)) 
    print(f"Rank {rank} [before all-gather]: input = {input}, 
          output = {output}", flush=True)
    dist.all_gather_into_tensor(output_tensor=output, 
                                input_tensor=input, async_op=False)
    print(f"Rank {rank} [after all-gather]: input = {input}, 
          output = {output}", flush=True)
    cleanup()
\end{minted}

Let's walk through some examples.

\begin{minted}{python}
def all_reduce(rank: int, world_size: int, num_elements: int):
    setup(rank, world_size)
    
    # Create tensor
    data = torch.randn(num_elements, device=cuda(rank))
    
    # Warmup
    dist.all_reduce(tensor=data, 
                    op=dist.ReduceOp.SUM, async_op=False)
    torch.cuda.synchronize()
    dist.barrier()
    # Perform all-reduce
    start_time = time.time()
    dist.all_reduce(tensor=data, 
                    op=dist.ReduceOp.SUM, 
                    async_op=False)
    torch.cuda.synchronize()
    dist.barrier()
    end_time = time.time()
    duration = end_time - start_time
    
    # Measure the effective bandwidth
    dist.barrier()
    size_bytes = data.element_size() * data.numel()
    sent_bytes = size_bytes * 2 * (world_size - 1)
    total_duration = world_size * duration
    bandwidth = sent_bytes / total_duration
    
    cleanup()  
\end{minted}

Notes: 

\begin{itemize}
    \item Effective bandwidth =  \(2 \times\)  size\_bytes / total\_duration.
    \item Independent of world\_size.
    \item Independent of topology (ring or tree).
\end{itemize}

\begin{minted}{python}
def reduce_scatter(rank: int, world_size: int, num_elements: int):
    setup(rank, world_size)
    
    # Create input and outputs
    input = torch.randn(world_size, num_elements, device=cuda(rank))
    output = torch.empty(num_elements, device=cuda(rank))
    
    # Warmup
    dist.reduce_scatter_tensor(
        input=input,output=output, 
        op=dist.ReduceOp.SUM, async_op=False)
    torch.cuda.synchronize() 
    dist.barrier()
    # Perform reduce-scatter
    start_time = time.time()
    dist.reduce_scatter_tensor(
        output=output, input=input, 
        op=dist.ReduceOp.SUM, async_op=False)
    torch.cuda.synchronize()
    dist.barrier()
    end_time = time.time()
    duration = end_time - start_time 
    
    # Measure the effective bandwidth
    dist.barrier()
    data_bytes = input.element_size() * input.numel()
    sent_bytes = data_bytes * (world_size - 1)
    total_duration = world_size * duration
    bandwidth = sent_bytes / total_duration
    cleanup()  
\end{minted}

Notes:

\begin{itemize}
    \item all-reduce \(=\) reduce-scatter \(+\) all-gather.
    \item all-reduce moves \(2 \times\) the data in \(2 \times\) the time compared to reduce-scatter.
\end{itemize}

\subsection*{Distributing the Training}

Walk through bare-bones implementations of each strategy on deep MLPs. Recall that MLPs are the compute bottleneck in Transformers, so this is representative. \\

What's missing?

\begin{itemize}
    \item Communication/computation overlap.
    \item More general models (with attention, etc).
    \item Other forms of parallelism (e.g., sequence parallelism, expert parallelism, combinations).
    \item Jax/TPUs: just define the model, the sharding strategy, and the Jax compiler handles the rest.
    \item But we're doing PyTorch so you can see how one builds up from the primitives.
\end{itemize}

\subsection*{Data Parallelism}

Sharding strategy: each rank gets a slice of the data.

\begin{itemize}
    \item Losses are different across ranks (computed on local data).
    \item Gradients are all-reduced to be the same across ranks.
    \item Therefore, parameters remain the same across ranks.
\end{itemize}

\begin{minted}{python}
def data_parallelism_main(rank: int, 
    world_size: int, data: tensor, 
    num_layers: int, num_steps: int):
    setup(rank, world_size)  
    
    # Get the slice of data for this rank
    # --- B0 ---
    # --- B1 ---
    # --- B2 ---
    # --- B3 ---
    batch_size = data.size(0)  
    num_dim = data.size(1)  
    local_batch_size = int_divide(batch_size, world_size)  
    start_index = rank * local_batch_size  
    end_index = start_index + local_batch_size  
    data = data[start_index:end_index].to(cuda(rank))
    
    # Create MLP parameters params[0], ..., params[num_layers - 1]
    params = [
        get_init_params(num_dim, num_dim, rank) 
        for layer in range(num_layers)]
    optimizer = torch.optim.AdamW(params, lr=1e-3)
    for step in range(num_steps):
        # Forward pass
        x = data
        for param in params:
            x = x @ param
            x = F.gelu(x)
        loss = x.square().mean()
        # Backward pass
        loss.backward()
        # Sync gradients across workers
        for param in params:
            dist.all_reduce(tensor=param.grad, 
                            op=dist.ReduceOp.AVG, async_op=False)            
        # Update parameters
        optimizer.step()
    cleanup()      
\end{minted}

\subsection*{Tensor Parallelism}

Sharding strategy: each rank gets part of each layer, transfer all data/activations.

\begin{minted}{python}
def tensor_parallelism_main(rank: int, world_size: int, 
    data: tensor, num_layers: int):
    setup(rank, world_size)  
    data = data.to(cuda(rank))
    batch_size = data.size(0)  
    num_dim = data.size(1)  
    local_num_dim = int_divide(num_dim, world_size) 
    
    # Create model
    #  |  |  |  |
    # W0 W1 W2 W3
    #  |  |  |  |
    params = [
        get_init_params(num_dim, local_num_dim, rank) 
        for layer in range(num_layers)]
    
    # Forward pass
    x = data
    for layer in range(num_layers):
        # Compute activations (batch_size x local_num_dim)
        x = x @ params[layer]
        x = F.gelu(x)
        # Allocate memory for activations
        activations = [
            torch.empty(batch_size, local_num_dim, 
            device=cuda_if_available(rank)) 
            for _ in range(world_size)]
        # Send activations via all gather
        dist.all_gather(
            tensor_list=activations, 
            tensor=x, async_op=False)
        # Concatenate them to get batch_size x num_dim
        x = torch.cat(activations, dim=1)
  
    # Backward pass: homework exercise
    cleanup()    
\end{minted}

\subsection*{Pipeline Parallelism}

Sharding strategy: each rank gets subset of layers, transfer all data/activations.

\begin{minted}{python}
def pipeline_parallelism_main(rank: int, world_size: int, 
    data: tensor, num_layers: int, num_micro_batches: int):
    setup(rank, world_size)  
    
    # Use all the data
    data = data.to(cuda_if_available(rank))
    batch_size = data.size(0)  
    num_dim = data.size(1)  
    
    # Split up layers
    local_num_layers = int_divide(num_layers, world_size)  
    # Each rank gets a subset of layers
    local_params = [
        get_init_params(num_dim, num_dim, rank) 
        for layer in range(local_num_layers)]  
    
    # Forward pass
    # Break up into micro batches to minimize the bubble
    micro_batch_size = int_divide(batch_size, num_micro_batches)  
    if rank == 0:
        # The data
        micro_batches = data.chunk(chunks=num_micro_batches, dim=0)
    else:
        # Allocate memory for activations
        micro_batches = [
            torch.empty(micro_batch_size, num_dim,
            device=cuda_if_available(rank)) 
            for _ in range(num_micro_batches)]
    for x in micro_batches:
        # Get activations from previous rank
        if rank - 1 >= 0:
            dist.recv(tensor=x, src=rank - 1)
        # Compute layers assigned to this rank
        for param in local_params:
            x = x @ param
            x = F.gelu(x)
        # Send to the next rank
        if rank + 1 < world_size:  
            dist.send(tensor=x, dst=rank + 1)

    # Backward pass: homework exercise
    cleanup()  
\end{minted}

\subsection*{Collective Parallelism}

\begin{itemize}
    \item Many ways to parallelize: data (batch), tensor/expert (width), pipeline (depth), sequence (length).
    \item Data parallelism: DDP (all-reduce), FSDP/ZeRO (all-gather + reduce-scatter).
    \item Tensor parallelism: requires very fast interconnects (e.g., NVLink).
    \item Pipeline parallelism: can work with slow interconnects, but need to work to reduce pipeline bubbles.
    \item Can re-compute or store in memory or store in another GPUs memory and communicate.
    \item Hardware is getting faster, but will always want bigger models, so will have this hierarchical structure.
\end{itemize}

\end{document}
